\documentclass[11pt]{article}
\pdfinfoomitdate=1
\pdftrailerid{}
\pdfsuppressptexinfo=15
\usepackage[T1]{fontenc}
\usepackage{lmodern,amsmath,amssymb,amsthm}
\usepackage[margin=1in]{geometry}
\usepackage{microtype,booktabs,needspace,multicol}
\usepackage[colorlinks=true,linkcolor=blue,urlcolor=blue]{hyperref}
\numberwithin{equation}{section}
\numberwithin{table}{section}
\title{Finite certificates for Hamiltonian thresholds\\in planar recolouring graphs}
\author{}
\date{}
\makeatletter
\renewcommand{\@maketitle}{%
 \newpage\null\vskip 1em
 \begin{center}
 {\LARGE \@title\par}
 \ifx\@author\@empty\else
 \vskip 1em{\large\begin{tabular}[t]{c}\@author\end{tabular}\par}
 \fi
 \ifx\@date\@empty\else\vskip .7em{\large\@date\par}\fi
 \end{center}\vskip 1em}
\makeatother
\begin{document}
\maketitle
\begin{abstract}
This supplement contains the literal local certificates used in
\emph{Hamiltonian and Hamilton-connected thresholds for planar recolouring graphs}.
The supplied Python verifiers read the displayed words and delivered
literal arrays directly and check the
complete finite case coverage. The inductions and arbitrary-order
restoration arguments are proved in the article.
\end{abstract}
Run \texttt{python3 verify\_all.py} to check the certificates and rejection
of corrupted data. The exact verifiers require only the Python standard
library; \texttt{README.md} gives their input contract and build commands.
\appendix
\section{Adjacent-pair endpoint certificates}
\label{app:pair-five-certificates}
For the article's fibre $F(A,B)$, write $ab$ for $(a,b)$ and use type
$(|A|,|B|,|A\cap B|)$. Each endpoint word visits every state once;
successive states differ in one coordinate. Table~\ref{tab:pair-five-counts}
includes list-preserving colour permutations, allowable coordinate
exchange and path reversal. Exact subset dynamic programming checks
all feasible and excluded endpoint pairs.

\begin{table}[htbp]
\centering
\caption{Adjacent-pair endpoint certificates and their symmetry coverage.}
\label{tab:pair-five-counts}
\begin{tabular}{crrr}
\toprule
Type & $|V(F)|$ & Words & Feasible unordered pairs\\
\midrule
 331&8&6&28\\
 332&7&7&19\\
 333&6&1&6\\
 342&10&16&45\\
 343&9&8&36\\
 443&13&11&78\\
 444&12&4&66\\
\bottomrule
\end{tabular}
\end{table}


\subsection{Edge-covering cycles}
The columns give the type, lists separated by a slash, and one cyclic
word. Its list-preserving symmetry orbit, including allowable coordinate
exchange, covers every edge.
% BEGIN PAIR-CYCLES
\begin{center}\small\ttfamily
\begin{tabular}{ccl}
331 & 012/034 & 03 04 14 10 20 24 23 13\\
332 & 012/013 & 23 03 01 21 20 10 13\\
333 & 012/012 & 10 20 21 01 02 12\\
342 & 012/0134 & 24 04 01 03 13 14 10 20 21 23\\
343 & 012/0123 & 12 10 20 21 01 02 03 23 13\\
443 & 0123/0124 & 32 12 10 20 21 01 02 04 24 14 34 31 30\\
444 & 0123/0123 & 23 21 20 10 12 02 32 30 31 01 03 13
\end{tabular}
\end{center}
% END PAIR-CYCLES

\Needspace{10\baselineskip}
\paragraph{The 53 endpoint words.}
The list pairs are given in the preceding cycle table.
\begin{multicols}{2}
\small
\noindent\begin{minipage}{\linewidth}
\noindent For $331$:
\begin{verbatim}
03 13 23 24 20 10 14 04
03 23 20 24 04 14 13 10
03 23 20 24 04 14 10 13
03 23 13 10 20 24 04 14
13 10 20 24 23 03 04 14
13 14 10 20 23 03 04 24
\end{verbatim}
\end{minipage}\par\medskip
\noindent\begin{minipage}{\linewidth}
\noindent For $332$:
\begin{verbatim}
01 21 23 20 10 13 03
01 21 20 23 03 13 10
01 03 23 21 20 10 13
01 21 20 10 13 03 23
03 01 21 23 20 10 13
03 01 21 23 13 10 20
03 01 21 20 10 13 23
\end{verbatim}
\end{minipage}\par\medskip
\noindent\begin{minipage}{\linewidth}
\noindent For $333$:
\begin{verbatim}
01 21 20 10 12 02
\end{verbatim}
\end{minipage}\par\medskip
\noindent\begin{minipage}{\linewidth}
\noindent For $342$:
\begin{verbatim}
01 21 24 20 23 13 10 14 04 03
01 21 23 20 24 14 04 03 13 10
01 21 24 23 20 10 14 04 03 13
01 21 23 24 14 04 03 13 10 20
01 04 14 24 23 03 13 10 20 21
01 21 24 20 10 13 14 04 03 23
03 23 24 14 13 10 20 21 01 04
03 23 24 20 21 01 04 14 10 13
03 13 10 20 24 23 21 01 04 14
03 23 24 21 01 04 14 13 10 20
03 13 23 24 20 10 14 04 01 21
03 01 21 20 24 04 14 10 13 23
03 23 13 14 10 20 21 01 04 24
20 24 23 13 10 14 04 03 01 21
20 10 13 14 04 24 21 01 03 23
23 21 20 10 14 13 03 01 04 24
\end{verbatim}
\end{minipage}\par\medskip
\noindent\begin{minipage}{\linewidth}
\noindent For $343$:
\begin{verbatim}
01 21 23 20 10 12 13 03 02
01 21 20 23 13 10 12 02 03
01 21 20 23 13 03 02 12 10
01 21 23 20 10 13 03 02 12
01 21 23 20 10 12 02 03 13
01 02 03 23 13 12 10 20 21
01 21 20 10 13 12 02 03 23
03 23 20 21 01 02 12 10 13
\end{verbatim}
\end{minipage}\par\medskip
\noindent\begin{minipage}{\linewidth}
\noindent For $443$:
\begin{verbatim}
01 31 34 24 21 20 30 32 12 10 14 04 02
01 21 31 34 32 30 20 24 14 10 12 02 04
01 21 31 34 32 30 20 24 14 04 02 12 10
01 31 32 30 34 24 21 20 10 14 04 02 12
01 31 32 30 34 24 21 20 10 12 02 04 14
01 21 20 30 32 31 34 14 10 12 02 04 24
01 21 31 32 30 20 24 14 10 12 02 04 34
04 24 34 32 31 30 20 21 01 02 12 10 14
04 34 14 24 20 21 01 31 32 02 12 10 30
04 34 32 30 20 21 24 14 10 12 02 01 31
04 24 20 30 32 31 21 01 02 12 10 14 34
\end{verbatim}
\end{minipage}\par\medskip
\noindent\begin{minipage}{\linewidth}
\noindent For $444$:
\begin{verbatim}
01 31 21 23 20 30 32 12 10 13 03 02
01 21 31 32 30 20 23 13 03 02 12 10
01 21 31 32 30 10 20 23 13 03 02 12
01 21 31 32 30 20 10 13 12 02 03 23
\end{verbatim}
\end{minipage}\par\medskip

\end{multicols}

\Needspace{8\baselineskip}
\section{Certificates for the \texorpdfstring{$(2,3,2)$}{(2,3,2)}-path restoration}
\label{app:path232-certificates}
The fibres $B_t$, allowed colour images and path reversals, and special
edge sets $P_t$ are defined in the main text. The following literal
cycles supply the selected-layer identities used in
the safe-common-edge lemma in the article.

\subsection{The three normalized cycles}
In the cyclic words below the coordinates are $uvw$. The rows for
$001,010,012$ specify $D_E,D_M,D_A$, respectively.

% BEGIN PATH232-CYCLES
\begin{equation}
\label{eq:path232-cycles}
\begin{array}{c|l}
001&120\ 123\ 323\ 313\ 213\ 210\ 230\ 130\ 132\ 232\ 212\ 312\ 310\ 320\\[2pt]
010&101\ 121\ 321\ 323\ 123\ 103\ 303\ 302\ 301\\
   &201\ 203\ 202\ 102\ 132\ 232\ 231\ 131\\[2pt]
012&131\ 101\ 103\ 203\ 303\ 301\ 201\ 231\\
   &230\ 130\ 120\ 123\ 323\ 320\ 321\ 121
\end{array}
\end{equation}
% END PATH232-CYCLES

\subsection{Selected-layer identities and capacities}
Order the end-, middle-, and all-distinct types as $E<M<A$.
For distinct $a,b,c,d$, let $E(aab),M_0(aba),A(abc)$ be the images of
$D_E,D_M,D_A$ under $(0,1,2,3)\mapsto(a,b,c,d)$, and let $M_1$ be
the reversal of $M_0$. Unused-colour interchange fixes $E,M_0,M_1$;
reversal followed by endpoint-colour interchange fixes $A$. Thus
$E,A$ each have one allowed edge set and $M$ has the two stated sets.
Up to these symmetries, the good transitions to a type no smaller are
\[
E(aab):aac,abb,acb,bab,cab;\quad
M(aba):aca,cba;\quad A(abc):abd,adc.
\]
These exhaust one-coordinate changes; bad targets and smaller types
are excluded, and unused colours and the two ends are identified
by the stated symmetries.

In $E(aab)$, the middle-$b$ layer is $2P_3$ and the middle-$c,d$
squares are $P_4$. In $M_0(aba)$, the middle-$a$ layer is $P_8$
with isolated $bab$, and the middle-$c,d$ squares are $P_4$;
reverse for $M_1$. In $A(abc)$, the middle-$a,c$ layers are $P_6$
and the middle-$d$ square is $P_3$ with isolated $bdb$.
These are subgraphs of the respective selected cycles.

Each cycle has four middle-changing edges. For $E$, middle colour $b$
joins $c,d$ twice each, with last coordinate $a$ twice and $c,d$
once each. For $M_0$, middle colour $a$ joins $c,d$ twice each and
every edge has first coordinate $b$; reverse for $M_1$. For $A$,
middle colour $d$ joins $a,c$ twice each, with first coordinate in
$\{b,c\}$ and last coordinate in $\{a,b\}$. These paths and counts
are read directly from~\eqref{eq:path232-cycles}.

For common endpoint lists $L,R$ and distinct allowed middle colours $v,w$,
the layer $H_v$ is
$K_{|L\setminus\{v\}|}\square K_{|R\setminus\{v\}|}$, and there are
$|L\setminus\{v,w\}|\,|R\setminus\{v,w\}|$ edges between middle
colours $v,w$. Removing the special edges gives the safe universe
size $m$; the original cyclic words give the retained counts $x,y$.
The following seven families supply the capacities used in the proof:
\[
\begin{array}{c|c|c}
t\to s&\text{blocks}&(m,x,y)\\\hline
aab\to aac&H_*&(4,3,3)\\
aab\to abb&H_c&(2,2,2)\\
aab\to acb&H_*&(2,2,2)\\
aab\to cab&H_c,H_d&(3,2,2),(1,1,1)\\
aba\to aca&H_d&(4,3,3)\\
abc\to abd&H_*&(6,4,4)\\
abc\to adc&H_a,H_c&(9,5,5)\text{ each}
\end{array}
\]
\Needspace{8\baselineskip}
\section{Literal pentagon cycles, path blocks and common-port seeds}
\label{app:pentagon-factor-words}
Five digits record the five colours in cyclic vertex order. In six-digit
words the first digit is the path layer. The two constant-type component
cycles are printed separately. Cycle words are cyclic; the
51-state and 47-state paths for $00112$ are linear.
A port is a pair of five-digit words.
The ordered path blocks use the specified cycles and the two
additional identities below; no eleven-type factor-word list is needed.
These literal identities are part of the proof, not solver verdicts.

\begin{table}[!htbp]
\centering\footnotesize
\caption{Ordered pairs for $00001$. Each pair $a_i b_i$ is a block;
the connector to the next block is $a_i a_{i+1}$ when $\epsilon_i=0$
and $b_i a_{i+1}$ when $\epsilon_i=1$.}
\label{tab:pentagon-matching-chain}
% BEGIN MATCHING-CHAIN-00001
\begingroup\ttfamily
\begin{tabular}{rllc@{\qquad}rllc}
$i$ & $a_i$ & $b_i$ & $\epsilon_i$ & $i$ & $a_i$ & $b_i$ & $\epsilon_i$\\
\hline
1 & 12132 & 32132 & 0 & 12 & 21310 & 31310 & 1\\
2 & 12130 & 32130 & 0 & 13 & 31312 & 32312 & 0\\
3 & 12120 & 32120 & 0 & 14 & 31212 & 31232 & 1\\
4 & 13120 & 23120 & 0 & 15 & 31230 & 21230 & 1\\
5 & 13123 & 23123 & 0 & 16 & 21210 & 31210 & 0\\
6 & 12123 & 12323 & 1 & 17 & 21213 & 23213 & 1\\
7 & 12313 & 12312 & 1 & 18 & 23210 & 13210 & 1\\
8 & 12310 & 32310 & 0 & 19 & 13212 & 13213 & 0\\
9 & 12320 & 32320 & 1 & 20 & 13232 & 13132 & 1\\
10 & 31320 & 21320 & 1 & 21 & 13130 & 23130 & 0\\
11 & 21323 & 21313 & 1 & 22 & 13230 & 23230 & --\\
\end{tabular}
\endgroup
% END MATCHING-CHAIN-00001
\end{table}

\paragraph{Path for 00112 (51 words).}
% BEGIN HP-00112
\begingroup\footnotesize\ttfamily\begin{flushleft}
12303 12323 12023 13023 13020 12020 12320 32320 31320 21320\\
21323 21303 21203 21201 21301 21321 31321 32321 32301 31301\\
31201 31231 31230 21230 21231 21031 21030 31030 31031 32031\\
32021 31021 31020 32020 32030 12030 13030 13230 23230 23030\\
23020 21020 21021 21023 23023 23021 23031 23231 23201 23203\\
13203\\
\end{flushleft}\endgroup
% END HP-00112

\paragraph{Constant-boundary cycles (30 words).}
% BEGIN FACTOR-00000
\begingroup\footnotesize\ttfamily\begin{flushleft}
12123 12323 12313 12312 32312 31312 31212 31232 31231 21231 23231 23131 23121 23123 13123\\
12132 13132 13232 13212 13213 23213 21213 21313 21323 21321 31321 32321 32121 32131 32132\\
\end{flushleft}\endgroup
% END FACTOR-00000
\paragraph{Cycle 00011 (42 words).}
% BEGIN C1-00011
\begingroup\small\ttfamily\begin{flushleft}
12102 13102 13103 23103 23123 13123 13120 23120 23130 13130\\
13230 23230 21230 31230 31232 31202 31302 32302 12302 12303\\
12103 12123 12323 12320 32320 31320 21320 21323 21303 21203\\
23203 13203 13202 13232 13132 12132 12130 12120 32120 32130\\
32132 32102\\
\end{flushleft}\endgroup
% END C1-00011
\paragraph{Cycle 00012 (47 words).}
% BEGIN C1-00012
\begingroup\small\ttfamily\begin{flushleft}
12103 12123 13123 13120 12120 32120 32130 12130 13130 13230\\
23230 21230 31230 31231 21231 23231 23201 23101 23131 23130\\
23120 23121 23123 23103 13103 13203 23203 21203 21201 31201\\
31301 32301 32101 32131 32121 32321 31321 21321 21301 21303\\
21323 21320 31320 32320 12320 12323 12303\\
\end{flushleft}\endgroup
% END C1-00012
\paragraph{Cycle 00101 (64 words).}
% BEGIN C1-00101
\begingroup\small\ttfamily\begin{flushleft}
12010 12020 32020 31020 21020 23020 13020 13010 23010 23210\\
13210 13212 13232 13032 13012 12012 32012 31012 31312 32312\\
12312 12310 32310 31310 21310 21320 31320 32320 12320 12323\\
12313 12013 12023 13023 23023 21023 21323 21313 21013 23013\\
13013 13213 23213 21213 21210 31210 31212 31232 31032 32032\\
12032 12030 32030 31030 31230 21230 23230 13230 13030 23030\\
21030 21010 31010 32010\\
\end{flushleft}\endgroup
% END C1-00101
\paragraph{Cycle 00102 (63 words).}
% BEGIN C1-00102
\begingroup\small\ttfamily\begin{flushleft}
12010 12030 32030 31030 31230 21230 23230 13230 13030 23030\\
21030 21031 23031 23231 21231 31231 31031 32031 32021 32321\\
31321 21321 21320 31320 32320 12320 12310 32310 31310 21310\\
21313 21323 21023 23023 13023 13020 13010 23010 21010 21020\\
23020 23021 21021 31021 31020 32020 12020 12023 12323 12313\\
12013 13013 23013 21013 21213 23213 13213 13210 23210 21210\\
31210 31010 32010\\
\end{flushleft}\endgroup
% END C1-00102
\paragraph{Cycle 00121 (58 words).}
% BEGIN C1-00121
\begingroup\small\ttfamily\begin{flushleft}
12010 13010 23010 23210 13210 13230 13030 12030 32030 31030\\
31032 32032 12032 13032 13232 13202 13212 13012 12012 32012\\
31012 31312 32312 12312 12310 32310 31310 21310 21210 21213\\
23213 13213 13203 23203 21203 21303 21313 21013 23013 13013\\
12013 12313 12303 12302 32302 31302 31202 31232 31212 31210\\
31230 21230 23230 23030 21030 21010 31010 32010\\
\end{flushleft}\endgroup
% END C1-00121
\paragraph{Cycle 00123 (57 words).}
% BEGIN C1-00123
\begingroup\small\ttfamily\begin{flushleft}
12010 12310 32310 31310 21310 21010 23010 13010 13210 13212\\
13202 13232 13032 13012 12012 32012 31012 31212 31232 31032\\
32032 12032 12030 32030 32031 31031 21031 21030 31030 31230\\
21230 21231 31231 31201 31202 31302 31312 32312 12312 12302\\
32302 32301 31301 21301 21201 23201 23231 23031 23030 13030\\
13230 23230 23210 21210 31210 31010 32010\\
\end{flushleft}\endgroup
% END C1-00123
\paragraph{Cycle 01012 (68 words).}
% BEGIN C1-01012
\begingroup\small\ttfamily\begin{flushleft}
10103 10203 13203 23203 20203 20303 10303 12303 12103 12123\\
10123 13123 23123 20123 20323 10323 12323 12320 32320 30320\\
10320 20320 20120 23120 13120 10120 30120 32120 12120 12130\\
32130 30130 10130 10230 30230 20230 23230 13230 13130 23130\\
20130 20131 23131 23231 20231 30231 30131 32131 32101 32301\\
30301 20301 20321 30321 32321 32121 30121 20121 23121 23101\\
23201 20201 30201 30101 20101 20103 23103 13103\\
\end{flushleft}\endgroup
% END C1-01012
\paragraph{Cycle 01023 (66 words).}
% BEGIN C1-01023
\begingroup\small\ttfamily\begin{flushleft}
10102 10302 30302 32302 12302 12312 10312 30312 32312 32310\\
12310 10310 30310 20310 20210 30210 10210 10212 30212 30202\\
10202 13202 13102 13132 13232 13212 13210 23210 23230 23130\\
23131 23231 23201 23101 20101 30101 30201 20201 20301 30301\\
32301 32101 32131 30131 30231 20231 20131 20130 20230 30230\\
10230 13230 13130 12130 32130 30130 10130 10132 10232 30232\\
30132 32132 12132 12102 32102 30102\\
\end{flushleft}\endgroup
% END C1-01023
\paragraph{C2 for 00112 (102 words).}
% BEGIN C2-00112
\begingroup\small\ttfamily\begin{flushleft}
012020 012030 032030 032031 032021 031021 021021 021031 031031\\
031231 031201 031301 032301 032321 032320 012320 012323 012303\\
112303 112323 112320 132320 132321 132301 131301 131201 131231\\
131230 031230 031030 021030 021230 121230 121030 131030 131020\\
121020 123020 123023 123021 123031 123030 123230 113230 113030\\
113020 113023 112023 012023 013023 013020 013030 013230 023230\\
023030 023031 023231 123231 121231 021231 021201 021301 021321\\
031321 131321 131021 131031 121031 121021 121023 121323 121303\\
021303 021203 121203 123203 113203 013203 023203 023201 123201\\
121201 121301 121321 121320 131320 031320 021320 021323 021023\\
023023 023021 023020 021020 031020 032020 132020 132021 132031\\
132030 112030 112020\\
\end{flushleft}\endgroup
% END C2-00112
\paragraph{Three-row insertion rule.}
Let $C_2$ be a Hamilton cycle of $A\square P_2$, with row digits
$0,1$. Suppose its
second row contains three disjoint edges $u_i v_i$, and the states
of $A$ partition into three paths with respective endpoints $u_i,v_i$.
Replace each cut edge $(1,u_i)(1,v_i)$ by its path in row $2$ and
the cross edges at both endpoints. These paths are internally
disjoint and cover the new row, so they subdivide three edges of
$C_2$ into one Hamilton cycle of $A\square P_3$.
If $D_A$ occurs in both old rows and in the new paths, and avoids
the three cut edges, it survives in every row of the derived seed.

\paragraph{Common ports 00112 (13 edges).}
% BEGIN PORTS-00112
\begingroup\small\ttfamily\begin{flushleft}
12023-13023 12320-32320 13020-13023 13030-13230\\
13230-23230 21020-23020 21020-31020 21030-21230\\
21030-31030 21301-21321 31201-31301 32030-32031\\
32301-32321\\
\end{flushleft}\endgroup
% END PORTS-00112
\paragraph{Generating the three-row seed for 00112.}
The following path $P$ contains forty-seven distinct states and
all thirteen ports above. Together with $12303,12323,13203,23203$,
its vertices are precisely the legal states of $00112$.
% BEGIN COREPATH-00112
\begingroup\small\ttfamily\begin{flushleft}
21201 21203 21303 21301 21321 21320 21323 21023 23023 23021\\
23031 23030 23230 13230 13030 12030 32030 32031 32021 32020\\
31020 21020 23020 13020 13023 12023 12020 12320 32320 31320\\
31321 32321 32301 31301 31201 31231 31230 21230 21030 31030\\
31031 31021 21021 21031 21231 23231 23201
\end{flushleft}\endgroup
% END COREPATH-00112
In row $1$ of the printed $C_2$, cut
\[
12303\text{--}12323,\qquad13203\text{--}23203,\qquad
21201\text{--}23201.
\]
The new row uses the first two edges as its two-point paths and
$P$ as the third path. All cut edges avoid $D_A$, which occurs
in both rows of $C_2$ and in $P$. The insertion rule therefore
gives the required three-row seed with all thirteen common ports.


\Needspace{8\baselineskip}
\section{Literal protected-pentagon certificates}
\label{app:protected-pentagon-words}
Boundary digit 4 denotes the full list. Colour words use only 0,1,2,3.
In six-digit cycle words the first digit is the path layer.
The first two pool edges for a Hamilton component join two or three
path rows; any two remaining pool edges supply block-chain reserves.
The protected set is the specified cycle edge set, or COMMON set
for a hard type, minus POOL. The earlier appendix supplies the
ordinary cycles, the two constant-type 15-cycles, and the 00112 seeds.
An interface record gives source type:component, position, new colour,
normalized target type:component, and the actual common edge.
Positions start at zero. Two records per line are separated by semicolons.
\paragraph{PC2 for 00001 (88 words).}
% BEGIN PC2-00001
\begingroup\footnotesize\ttfamily\begin{flushleft}
012120 012130 112130 112120 132120 132130 132132 112132 113132\\
113232 113230 113210 123210 123213 113213 113212 013212 013213\\
023213 021213 021210 023210 013210 013230 013232 013132 012132\\
032132 032130 032120 032320 032310 012310 012320 012323 012313\\
012312 032312 031312 031310 031210 031212 031232 031230 021230\\
023230 123230 121230 131230 131232 131212 131210 121210 121213\\
121313 121323 021323 021313 021310 121310 121320 021320 031320\\
131320 132320 112320 112310 132310 131310 131312 132312 112312\\
112313 112323 112123 012123 013123 023123 123123 113123 113120\\
123120 123130 113130 013130 023130 023120 013120\\
\end{flushleft}\endgroup
% END PC2-00001
\paragraph{Common ports 00001 (13 edges).}
% BEGIN COMMON-00001
\begingroup\footnotesize\ttfamily\begin{flushleft}
12132-13132 12132-32132 12310-32310 13120-23120\\
13123-23123 13210-23210 13230-13232 21210-21213\\
21230-23230 21230-31230 21313-21323 31212-31232\\
32130-32132\\
\end{flushleft}\endgroup
% END COMMON-00001
\paragraph{Generating the three-row seed for 00001.}
The following path $P$ contains forty distinct states and all
thirteen ports above. Together with $12312,12313,13212,13213$,
its vertices are precisely the legal states of $00001$.
% BEGIN COREPATH-00001
\begingroup\footnotesize\ttfamily\begin{flushleft}
13210 23210 23213 21213 21210 21310 21313 21323 21320 31320\\
32320 32120 12120 13120 23120 23123 13123 12123 12323 12320\\
12310 32310 32312 31312 31310 31210 31212 31232 31230 21230\\
23230 23130 13130 12130 32130 32132 12132 13132 13232 13230\\
\end{flushleft}\endgroup
% END COREPATH-00001
In row $1$ of the printed $C_2$, cut
\[
12312\text{--}12313,\qquad13212\text{--}13213,\qquad
13210\text{--}13230.
\]
The new row uses the first two edges and $P$. All cuts avoid
$D_A$, which occurs in both rows of $C_2$ and in $P$. Apply the
insertion rule in Appendix~\ref{app:pentagon-factor-words} to obtain
the three-row seed, retaining every listed common port.
\paragraph{Cycle for 00004 (54 words).}
% BEGIN FHC-00004
\begingroup\footnotesize\ttfamily\begin{flushleft}
12120 12130 12132 13132 13130 23130 23230 23210 13210 13230\\
13232 13212 13213 23213 21213 21210 31210 31212 31232 31231\\
31230 21230 21231 23231 23131 23121 23123 23120 13120 13123\\
12123 12323 12313 12312 12310 32310 32312 31312 31310 21310\\
21313 21323 21321 21320 31320 31321 32321 32121 32131 32132\\
32130 32120 32320 12320\\
\end{flushleft}\endgroup
% END FHC-00004
\paragraph{Cycle for 00014 (60 words).}
% BEGIN FHC-00014
\begingroup\footnotesize\ttfamily\begin{flushleft}
12102 13102 13132 13232 13202 13203 13103 12103 12123 13123\\
13120 12120 32120 32320 12320 12323 12303 12302 32302 31302\\
31202 31232 31230 21230 23230 13230 13130 23130 23120 23121\\
23123 23103 23203 23201 23101 23131 23231 21231 31231 31201\\
21201 21203 21303 21323 21320 31320 31321 21321 21301 31301\\
32301 32321 32121 32101 32131 32130 12130 12132 32132 32102\\
\end{flushleft}\endgroup
% END FHC-00014
\paragraph{Cycle for 00104 (78 words).}
% BEGIN FHC-00104
\begingroup\footnotesize\ttfamily\begin{flushleft}
12010 12030 32030 31030 31230 21230 23230 13230 13030 23030\\
21030 21031 23031 23231 21231 31231 31031 32031 32021 32321\\
32320 12320 12323 12313 12013 12023 13023 23023 21023 21020\\
31020 32020 12020 13020 23020 23021 21021 31021 31321 31320\\
21320 21321 21323 21313 21310 31310 32310 12310 12312 32312\\
31312 31012 32012 12012 13012 13013 23013 21013 21213 23213\\
13213 13212 13232 13032 12032 32032 31032 31232 31212 31210\\
21210 23210 13210 13010 23010 21010 31010 32010\\
\end{flushleft}\endgroup
% END FHC-00104
\paragraph{Cycle for 00114 (62 words).}
% BEGIN FHC-00114
\begingroup\footnotesize\ttfamily\begin{flushleft}
12020 12320 32320 31320 21320 21323 21023 23023 23021 23031\\
21031 21021 21020 21030 21230 31230 31030 32030 12030 13030\\
13230 23230 23030 23020 13020 13023 12023 12323 12303 12302\\
32302 31302 31202 31232 31032 32032 12032 13032 13232 13202\\
13203 23203 21203 21303 21301 21321 31321 32321 32301 31301\\
31201 21201 23201 23231 21231 31231 31031 32031 32021 31021\\
31020 32020\\
\end{flushleft}\endgroup
% END FHC-00114
\paragraph{Cycle for 00124 (71 words).}
% BEGIN FHC-00124
\begingroup\footnotesize\ttfamily\begin{flushleft}
12010 12030 32030 32031 31031 31231 21231 21031 21030 31030\\
31230 21230 23230 13230 13030 23030 23031 23231 23201 23203\\
21203 21303 21301 21201 31201 31202 31302 31301 32301 32302\\
12302 12303 12313 12013 13013 13010 13210 13212 13012 12012\\
32012 31012 31312 32312 12312 12310 32310 31310 31210 31212\\
31232 31032 32032 12032 13032 13232 13202 13203 13213 23213\\
21213 21313 21013 23013 23010 23210 21210 21310 21010 31010\\
32010\\
\end{flushleft}\endgroup
% END FHC-00124
\paragraph{Cycle for 01014 (88 words).}
% BEGIN FHC-01014
\begingroup\footnotesize\ttfamily\begin{flushleft}
10102 10132 13132 13232 10232 30232 30132 32132 12132 12130\\
32130 30130 10130 10230 30230 20230 23230 13230 13130 23130\\
20130 20131 23131 23231 20231 30231 30131 32131 32101 32301\\
30301 20301 20101 30101 30201 20201 23201 23101 23121 20121\\
30121 32121 32321 30321 20321 20323 10323 12323 12320 32320\\
30320 10320 20320 20120 23120 13120 12120 32120 30120 10120\\
10123 12123 13123 23123 20123 20103 20203 20303 10303 10302\\
30302 32302 12302 12303 12103 10103 10203 13203 23203 23103\\
13103 13102 13202 10202 30202 30102 32102 12102\\
\end{flushleft}\endgroup
% END FHC-01014
\paragraph{Cycle for 01024 (85 words).}
% BEGIN FHC-01024
\begingroup\footnotesize\ttfamily\begin{flushleft}
10102 10202 30202 30212 10212 13212 13210 23210 20210 30210\\
10210 10213 10203 20203 23203 13203 13213 23213 20213 20313\\
10313 12313 12312 12310 32310 30310 20310 10310 10312 30312\\
32312 32302 30302 10302 10303 20303 20103 23103 13103 10103\\
12103 12303 12302 12102 32102 30102 30101 30201 20201 23201\\
23101 20101 20301 30301 32301 32101 32131 30131 30231 20231\\
23231 23131 20131 20130 23130 13130 13230 23230 20230 30230\\
10230 10130 30130 32130 12130 12132 32132 30132 30232 10232\\
10132 13132 13232 13202 13102\\
\end{flushleft}\endgroup
% END FHC-01024
\paragraph{Cycle for 01104 (74 words).}
% BEGIN FHC-01104
\begingroup\footnotesize\ttfamily\begin{flushleft}
10210 10310 20310 30310 32310 12310 12312 10312 30312 32312\\
32012 32032 32031 32021 32321 30321 20321 20320 10320 30320\\
32320 12320 12020 12010 32010 32020 32030 12030 12032 13032\\
13232 13212 13012 12012 12013 12313 10313 20313 20323 10323\\
12323 12023 13023 23023 23013 13013 13213 23213 20213 10213\\
10212 10232 10230 13230 23230 23030 13030 13020 23020 23021\\
23031 23231 20231 30231 30232 30212 30210 30230 20230 20210\\
23210 23010 13010 13210\\
\end{flushleft}\endgroup
% END FHC-01104
\paragraph{Cycle for 01124 (74 words).}
% BEGIN FHC-01124
\begingroup\footnotesize\ttfamily\begin{flushleft}
10202 10203 20203 23203 23201 20201 20301 20303 10303 12303\\
12302 32302 32301 30301 30201 30231 20231 20230 30230 10230\\
13230 23230 23231 23031 23030 13030 12030 32030 32031 32032\\
32012 32010 12010 12012 12032 13032 13232 10232 30232 30202\\
30212 10212 10213 10210 30210 20210 20310 30310 10310 12310\\
32310 32312 12312 12313 12013 13013 23013 23010 23210 13210\\
13010 13012 13212 13202 13203 13213 23213 20213 20313 10313\\
10312 30312 30302 10302\\
\end{flushleft}\endgroup
% END FHC-01124
\paragraph{Cycle for 01204 (83 words).}
% BEGIN FHC-01204
\begingroup\footnotesize\ttfamily\begin{flushleft}
10120 20120 23120 13120 13130 23130 20130 10130 30130 32130\\
12130 12132 32132 30132 10132 13132 13032 12032 32032 32012\\
12012 13012 13010 12010 32010 32020 12020 13020 23020 23010\\
23030 13030 12030 32030 32031 32131 30131 20131 23131 23031\\
23021 23121 20121 30121 32121 32021 32321 30321 20321 20320\\
10320 30320 32320 12320 12323 12023 13023 23023 23013 13013\\
12013 12313 12310 32310 32312 12312 10312 30312 30310 10310\\
20310 20313 10313 10323 20323 20123 23123 13123 10123 12123\\
12120 32120 30120\\
\end{flushleft}\endgroup
% END FHC-01204
\paragraph{Cycle for 01234 (82 words).}
% BEGIN FHC-01234
\begingroup\footnotesize\ttfamily\begin{flushleft}
10102 10302 30302 32302 12302 12303 10303 10323 20323 20303\\
20103 10103 12103 12123 13123 23123 20123 10123 10120 30120\\
32120 12120 13120 23120 20120 20320 10320 30320 32320 12320\\
12310 32310 30310 10310 20310 20313 10313 10312 30312 32312\\
12312 12313 12323 12023 13023 23023 23013 23010 13010 12010\\
12013 13013 13012 12012 32012 32010 32020 12020 13020 23020\\
23021 23121 20121 30121 32121 32021 32321 30321 20321 20301\\
30301 32301 32101 30101 20101 23101 23103 13103 13102 12102\\
32102 30102\\
\end{flushleft}\endgroup
% END FHC-01234
\paragraph{Pool for 00000:0.}
% BEGIN POOL-00000-0
\begingroup\footnotesize\ttfamily\begin{flushleft}
12123-13123 12312-12313 21231-31231 23121-23131 31212-31312\\
\end{flushleft}\endgroup
% END POOL-00000-0
\paragraph{Pool for 00000:1.}
% BEGIN POOL-00000-1
\begingroup\footnotesize\ttfamily\begin{flushleft}
12132-13132 13212-13213 21213-21313 21321-31321 32121-32131\\
\end{flushleft}\endgroup
% END POOL-00000-1
\begingroup
\setlength{\columnsep}{14pt}
\renewcommand{\paragraph}[1]{\par\addvspace{1ex}\noindent
  {\normalfont\small\bfseries #1}\par\nobreak}
\renewenvironment{flushleft}{\par\raggedright}{\par}
\begin{multicols}{2}
\paragraph{Pool for 00001:0.}
% BEGIN POOL-00001-0
\begingroup\footnotesize\ttfamily\begin{flushleft}
21210-21213 21230-31230\\
\end{flushleft}\endgroup
% END POOL-00001-0
\paragraph{Pool for 00011:0.}
% BEGIN POOL-00011-0
\begingroup\footnotesize\ttfamily\begin{flushleft}
12102-13102 13102-13103 13123-23123 13130-13230\\
13202-13203 21203-21303 31202-31302 32130-32132\\
\end{flushleft}\endgroup
% END POOL-00011-0
\paragraph{Pool for 00012:0.}
% BEGIN POOL-00012-0
\begingroup\footnotesize\ttfamily\begin{flushleft}
13130-13230 21301-21321 23120-23130 23121-23123\\
31201-31301 32101-32301 32121-32131 32121-32321\\
\end{flushleft}\endgroup
% END POOL-00012-0
\paragraph{Pool for 00101:0.}
% BEGIN POOL-00101-0
\begingroup\footnotesize\ttfamily\begin{flushleft}
12010-12020 12012-13012 12013-12023 12320-12323\\
13013-23013 13210-13212 21010-31010 23010-23210\\
\end{flushleft}\endgroup
% END POOL-00101-0
\paragraph{Pool for 00102:0.}
% BEGIN POOL-00102-0
\begingroup\footnotesize\ttfamily\begin{flushleft}
12013-13013 12313-12323 13013-23013 21013-21213\\
21023-21323 21031-23031 23020-23021 31031-31231\\
\end{flushleft}\endgroup
% END POOL-00102-0
\paragraph{Pool for 00112:0.}
% BEGIN POOL-00112-0
\begingroup\footnotesize\ttfamily\begin{flushleft}
21020-23020 21030-31030\\
\end{flushleft}\endgroup
% END POOL-00112-0
\paragraph{Pool for 00121:0.}
% BEGIN POOL-00121-0
\begingroup\footnotesize\ttfamily\begin{flushleft}
12010-13010 21010-21030 21010-31010 21210-21213\\
23010-23210 23030-23230 31012-31312 31030-32030\\
\end{flushleft}\endgroup
% END POOL-00121-0
\paragraph{Pool for 00123:0.}
% BEGIN POOL-00123-0
\begingroup\footnotesize\ttfamily\begin{flushleft}
13030-23030 13230-23230 21010-21310 21030-31030\\
21230-21231 31010-31210 31012-32012 31030-31230\\
\end{flushleft}\endgroup
% END POOL-00123-0
\paragraph{Pool for 01012:0.}
% BEGIN POOL-01012-0
\begingroup\footnotesize\ttfamily\begin{flushleft}
10103-13103 10123-13123 10130-30130 13130-13230\\
20101-30101 20230-30230 20301-20321 30131-30231\\
\end{flushleft}\endgroup
% END POOL-01012-0
\paragraph{Pool for 01023:0.}
% BEGIN POOL-01023-0
\begingroup\footnotesize\ttfamily\begin{flushleft}
12130-13130 20101-30101 20201-20301 20210-30210\\
20231-30231 23130-23131 30102-32102 30202-30212\\
\end{flushleft}\endgroup
% END POOL-01023-0
\paragraph{Pool for 00004:0.}
% BEGIN POOL-00004-0
\begingroup\footnotesize\ttfamily\begin{flushleft}
12312-12313 13130-13132 13212-13213 21210-21213\\
21321-21323 23121-23123 31231-31232 31310-31312\\
\end{flushleft}\endgroup
% END POOL-00004-0
\paragraph{Pool for 00014:0.}
% BEGIN POOL-00014-0
\begingroup\footnotesize\ttfamily\begin{flushleft}
12102-32102 12103-13103 12120-32120 12130-12132\\
13130-23130 21201-21203 23121-23123 32101-32121\\
\end{flushleft}\endgroup
% END POOL-00014-0
\paragraph{Pool for 00104:0.}
% BEGIN POOL-00104-0
\begingroup\footnotesize\ttfamily\begin{flushleft}
13012-13013 13210-23210 13212-13213 21020-21023\\
21320-21321 21321-21323 31320-31321 32320-32321\\
\end{flushleft}\endgroup
% END POOL-00104-0
\paragraph{Pool for 00114:0.}
% BEGIN POOL-00114-0
\begingroup\footnotesize\ttfamily\begin{flushleft}
12030-13030 21020-21021 21301-21321 21320-21323\\
23021-23023 31020-31021 31020-32020 32021-32031\\
\end{flushleft}\endgroup
% END POOL-00114-0
\paragraph{Pool for 00124:0.}
% BEGIN POOL-00124-0
\begingroup\footnotesize\ttfamily\begin{flushleft}
12310-32310 13210-13212 21030-31030 21201-31201\\
21213-21313 23010-23210 31010-32010 31301-31302\\
\end{flushleft}\endgroup
% END POOL-00124-0
\paragraph{Pool for 01014:0.}
% BEGIN POOL-01014-0
\begingroup\footnotesize\ttfamily\begin{flushleft}
10120-10123 10202-13202 10302-10303 12120-13120\\
12123-13123 13102-13103 20130-20131 20321-20323\\
\end{flushleft}\endgroup
% END POOL-01014-0
\paragraph{Pool for 01024:0.}
% BEGIN POOL-01024-0
\begingroup\footnotesize\ttfamily\begin{flushleft}
10103-13103 10132-13132 10302-10303 10312-30312\\
12130-12132 12312-12313 20103-20303 20213-20313\\
\end{flushleft}\endgroup
% END POOL-01024-0
\paragraph{Pool for 01104:0.}
% BEGIN POOL-01104-0
\begingroup\footnotesize\ttfamily\begin{flushleft}
10210-10310 12010-12020 12012-12013 12310-12312\\
13013-13213 20210-23210 23030-23230 32012-32312\\
\end{flushleft}\endgroup
% END POOL-01104-0
\paragraph{Pool for 01124:0.}
% BEGIN POOL-01124-0
\begingroup\footnotesize\ttfamily\begin{flushleft}
10202-10203 10210-10213 10230-13230 10310-30310\\
13010-13210 20210-20310 23230-23231 32310-32312\\
\end{flushleft}\endgroup
% END POOL-01124-0
\paragraph{Pool for 01204:0.}
% BEGIN POOL-01204-0
\begingroup\footnotesize\ttfamily\begin{flushleft}
10123-13123 12010-32010 12120-12123 12310-12313\\
12320-12323 23010-23020 30310-30312 32030-32031\\
\end{flushleft}\endgroup
% END POOL-01204-0
\paragraph{Pool for 01234:0.}
% BEGIN POOL-01234-0
\begingroup\footnotesize\ttfamily\begin{flushleft}
10102-30102 12102-32102 12123-13123 12302-32302\\
12310-12320 13012-13013 13102-13103 20101-30101\\
\end{flushleft}\endgroup
% END POOL-01234-0
\end{multicols}
\endgroup
% BEGIN CONSTANT-INTERFACE-FAMILIES
\paragraph{The thirty constant-source interfaces.}
Write $K_i=E(C_i)\setminus O_i$ for the two constant components.
Coordinate rotation fixes both components, since its order is five.
Each five-edge pool $O_i$ is a rotation orbit, so $K_i$ is invariant
under $\rho_r(x)_k=x_{k+r\pmod5}$. The following edges belong to
the unchanged protected set $K_H$ of type $00001$:
\[
a=31212\!-31232,\quad b=12132\!-32132,\quad
c=21313\!-21323,\quad d=13123\!-23123.
\]
When boundary entry $j$ changes from $0$ to $q\in\{1,2,3\}$,
its canonical normalizer is one of
\[
 N_{j,\epsilon}(x)_k=p_q(x_{j-4\epsilon+\epsilon k}),
 \qquad \epsilon\in\{+1,-1\},
\]
where $p_q$ is the first lexicographic permutation with $p_q(0)=0$
and $p_q(q)=1$. Indeed, the four repeated entries map to $0$ and
the unique changed entry maps to $1$ at position $4$.
\par\needspace{16\baselineskip}
The table gives $e_{i,q}\in K_i$ and its two normalized images;
these memberships follow directly from the printed words and pools.
\[
\begin{array}{cc|c|cc}
i&q&e_{i,q}&N_{0,+1}(e_{i,q})&N_{0,-1}(e_{i,q})\\ \hline
0&1&23121\!-23123&a&b\\
0&2&31231\!-31232&c&d\\
0&3&21231\!-23231&d&c\\
1&1&21213\!-23213&b&a\\
1&2&31321\!-32321&d&c\\
1&3&21321\!-21323&c&d
\end{array}
\]
At any position $j$, use $\rho_{-j}(e_{i,q})\in K_i$. Its image is
\[
 N_{j,\epsilon}(\rho_{-j}(e_{i,q}))
      =N_{0,\epsilon}(e_{i,q})\in K_H,
\]
so it lies in both actual protected sets. These six families cover
all thirty interfaces.
% END CONSTANT-INTERFACE-FAMILIES

\paragraph{Fixed-side source selectors (236 interfaces).}
Let $K$ be either fixed protected set and $B$ the opposite component,
with pool $P_B$ and a Hamilton cycle. At a common state $x$, put
$d=\deg_B(x)$ and $g=|(K\cap\delta_B(x))\setminus P_B|$.
If $d-g\le1$, every Hamilton cycle in $B$ takes two incident edges
at $x$, at least one in $K\setminus P_B$. Only $B$ needs a Hamilton cycle.

\[\begin{array}{c|l}t&S_t\\\hline
00001&13123,\ 21230,\ 23120,\ 32132\\
00004&12130,\ 13123,\ 32312\\
00011&12323,\ 21230,\ 21320\\
00012&12303,\ 13203,\ 13230,\ 31230\\
00014&12303,\ 31202,\ 31230\\
00101&12312,\ 13230,\ 13232,\ 21323,\ 31312\\
00102&13023,\ 13210,\ 13230,\ 23231,\ 31231,\ 32031,\ 32321\\
00104&12032,\ 12323,\ 13232,\ 23021,\ 23031\\
00112&12023,\ 13230,\ 31201\\
00114&12032,\ 21303\\
00121&12302,\ 12303,\ 21203,\ 21303\\
00123&12302,\ 12310,\ 23201,\ 32031\\
00124&12303,\ 23031,\ 32301\\
01012&10203,\ 13203,\ 23231,\ 32301\\
01014&12302,\ 23201,\ 23231,\ 32301\\
01023&12310,\ 13202,\ 20310,\ 23101\\
01024&23213,\ 32301,\ 32310\\
01104&20231,\ 30321,\ 32021\\
01124&12303,\ 23031,\ 32301\\
01204&10132,\ 13032,\ 23013,\ 30312,\ 32312\\
01234&23013,\ 23101\\
\end{array}\]


For each of the 236 nonconstant-source interfaces not listed individually
below, some $x\in S_t$ lies in both actual components and satisfies
$d-g\le1$, with fixed source $K$ against the restored target component,
or restored target $K$ against the source component. The opposite side
has its displayed Hamilton cycle; the printed protected sets and pools
use the original interface's first normalizer. Its degree is
\[
L_i(b,x)=\{0,1,2,3\}\setminus\{x_{i-1},x_{i+1},b_i\},\qquad
d=\sum_i(|L_i(b,x)|-1),
\]
where $b$ is its boundary in the original coordinates, omitting $b_i$
when starred. This proves the specified-template identities in all 236 cases.

\paragraph{The 47 remaining individual nonconstant-source interfaces.}
% BEGIN INTERFACES-ALL
\begingroup\scriptsize\ttfamily\begin{flushleft}
00001:0 1 2 00102:0 21313-21323 ; 00001:0 1 3 00102:0 21313-21323\\
00001:0 2 1 00101:0 12310-32310 ; 00001:0 4 2 00001:0 13120-23120\\
00001:0 4 3 00001:0 12310-32310 ; 00011:0 0 2 00112:0 12130-12132\\
00011:0 0 3 00112:0 12120-12130 ; 00101:0 0 1 00101:0 21020-23020\\
00101:0 1 1 00101:0 12013-12313 ; 00101:0 2 0 00001:0 12310-32310\\
00101:0 3 2 00121:0 12010-32010 ; 00101:0 3 3 00121:0 12010-32010\\
00101:0 4 0 00001:0 12032-32032 ; 00102:0 0 1 01012:0 21013-23013\\
00112:0 0 1 00102:0 21030-21230 ; 00112:0 0 2 00112:0 12023-13023\\
00112:0 1 3 00123:0 12320-32320 ; 00112:0 2 2 00121:0 12023-13023\\
00112:0 2 3 00123:0 13020-13023 ; 00112:0 3 2 00112:0 13030-13230\\
00112:0 4 3 00112:0 12320-32320 ; 00121:0 2 2 00112:0 12010-32010\\
00121:0 3 0 00101:0 12010-32010 ; 00121:0 4 2 00112:0 12013-13013\\
00123:0 2 2 00112:0 12030-32030 ; 01012:0 0 1 00102:0 20103-23103\\
01012:0 1 2 01012:0 10120-13120 ; 00104:0 0 1 01014:0 21010-31010\\
00104:0 1 2 01204:0 13010-23010 ; 00104:0 1 3 01204:0 12010-32010\\
00124:0 2 3 00124:0 12010-32010 ; 01014:0 0 1 00104:0 20101-30101\\
01014:0 0 2 01024:0 10102-12102 ; 01014:0 0 3 01024:0 10103-10203\\
01014:0 1 2 01024:0 10120-30120 ; 01014:0 1 3 01024:0 10102-12102\\
01014:0 2 2 01024:0 10102-10132 ; 01014:0 2 3 01024:0 10102-10132\\
01024:0 0 3 01234:0 10103-12103 ; 01024:0 1 2 01014:0 10102-13102\\
01024:0 2 3 01234:0 10102-13102 ; 01204:0 1 0 00104:0 12012-13012\\
01204:0 1 3 01204:0 10120-30120 ; 01234:0 0 2 01024:0 10102-10302\\
01234:0 1 3 01024:0 10102-10302 ; 01234:0 2 0 01024:0 10103-12103\\
01234:0 3 1 01024:0 10102-10302\\
\end{flushleft}\endgroup
% END INTERFACES-ALL

\clearpage
\section{Shared-hub path certificates}
\label{app:shared-hub-path-certificates}
Let $v_0,\ldots,v_{r-1}$ be the restored path, where $r\in\{4,5\}$.
Its endpoints have the same outside neighbour $w$; the other outside
neighbours of its internal vertices are distinct and different from $w$.
In the physical port order $t=(a,b_1,\ldots,b_{r-2})$, the fibre consists
of the proper path words with coordinate exclusions
\[
 (a,b_1,\ldots,b_{r-2},a).
\]
All colours in this appendix and its arrays lie in $\{0,1,2,3\}$.
The files \texttt{data/path4-family.json} and
\texttt{data/path5-family.json} supply literal cyclic arrays, one for
each normalized boundary. Normalize colours in first-appearance order,
then append unused colours in ascending order. Restore every cycle by
the inverse palette map, without reflecting the path or choosing
an additional relative map.

\begin{table}[htbp]
\centering
\caption{Complete shared-hub path coverage. Interfaces are directed.}
\label{tab:shared-hub-path-counts}
% BEGIN SHARED-HUB-PATH-COUNTS
\begin{tabular}{rrrrrr}
\toprule
Order & Templates & States & Boundaries & Interfaces & Min. edges\\
\midrule
4&5&160&64&576&3\\
5&15&1112&256&3072&5\\
\bottomrule
\end{tabular}
% END SHARED-HUB-PATH-COUNTS
\end{table}

Here ``States'' sums the lengths of the normalized arrays. The exact
reader reconstructs each legal fibre independently, then checks complete
coverage, no repeated state, and a one-coordinate change at every cycle
step, including the closing step. Raw array digits are checked before
palette-map indexing.

The reader expands to every labelled boundary and intersects the
\emph{chosen cycle edge sets} at every physical one-port change.
A change at $w$ changes both endpoint exclusions simultaneously; it is
one physical interface, included among the 576 or 3072 directed cases.
The endpoints of every shared edge are legal on both sides, so both
matching outside-recolouring edges exist. The displayed minima exceed
the two edges needed to select an outgoing splice after one incoming
edge has been removed.

The reader \texttt{verify\_shared\_hub\_paths.py} checks these finite
claims and records the size and SHA-256 of each input array.
The all-readers command also rejects missing or duplicate templates,
repeated states, invalid palette digits and inconsistent coverage counts.
Splicing along an
arbitrary bottom Hamilton path and the article's structural induction
remain analytic arguments; neither is inferred from finite replay.

\clearpage
\section{Degree-two configuration families}
\label{app:degree-two-family-certificates}
The seven files \texttt{data/degree-two/}\emph{kind}\texttt{-family.json}
give literal cycles for the following fixed internal models. Paths have
vertices $0,\ldots,n-1$ in order; the star has centre $0$.
Write $T_m$ for the cycle $0,\ldots,m-1,0$ with the additional vertex
$m$ and edge $0m$. A boundary is $t\in\{0,1,2,3\}^q$.
At vertex $v$, exclude the colours $t_i$ for $i\in I_v$.
In the tuples below a single index denotes its singleton set.
An empty set excludes no colour; a two-element set can exclude one
or two colours, according to their physical equality.
\[
\begin{array}{lcl}
\text{Kind}&X&(I_0,\ldots,I_{n-1})\\ \hline
\mathtt{star3}&K_{1,3}&(0,1,2,3)\\
\mathtt{path4heavy}&P_4&(0,\{1,2\},3,0)\\
\mathtt{path5heavynear}&P_5&(0,\{1,2\},3,4,0)\\
\mathtt{path5heavycenter}&P_5&(0,1,\{2,3\},4,0)\\
\mathtt{path6}&P_6&(0,1,2,3,4,0)\\
\mathtt{tailed5}&T_5&(0,\varnothing,1,2,3,4)\\
\mathtt{tailed6}&T_6&(0,\varnothing,1,2,3,4,5)
\end{array}
\]
The port counts $q$ are respectively $4,4,5,5,5,5,6$.
In every path model, port $0$ is the common outside neighbour of both
endpoints, so its move changes both endpoint exclusions at once.

Two port roles can represent one outside vertex only if every pair
of their incident internal vertices has distance at least three;
otherwise a triangle or a $4$-cycle results. The complete additional
alias pairs are
\[
\begin{array}{ll}
\mathtt{path6}:&(1,4),\\
\mathtt{tailed5}:&(1,4),\ (2,4),\\
\mathtt{tailed6}:&(0,2),\ (1,4),\ (1,5),\ (2,5),\ (3,5).
\end{array}
\]
The first four models have none. Each compatibility graph is
triangle-free, so no outside vertex occupies three roles. Several
disjoint pairs may be identified. A move changes one singleton role
or one listed pair, with equal old entries in the latter case.

\begin{table}[htbp]
\centering\small
\caption{Degree-two family coverage; interfaces are directed.}
\label{tab:degree-two-family-counts}
% BEGIN DEGREE-TWO-FAMILY-COUNTS
\begin{tabular}{lrrrrr}
\toprule
Kind & Templates & States & Boundaries & Interfaces & Min. edges\\
\midrule
star3&15&499&256&3072&4\\
path4heavy&15&376&256&3072&2\\
path5heavynear&51&2915&1024&15360&2\\
path5heavycenter&51&2921&1024&15360&2\\
path6&51&8593&1024&16128&9\\
tailed5&51&8535&1024&16896&12\\
tailed6&187&72197&4096&89088&19\\
\bottomrule
\end{tabular}
% END DEGREE-TWO-FAMILY-COUNTS
\end{table}

The primary reader reconstructs every legal fibre from the fixed models,
checks the literal cycles and verifies alias completeness. Every labelled
singleton or paired transition is enumerated; repeated relative-colour
intersections are cached. Palette normalization is as in
Appendix~\ref{app:shared-hub-path-certificates}. The separate direct reader
restores cycles to physical labels and intersects packed edge sets,
without importing the primary reader.

Both readers validate raw digits before indexing, enforce all seven file
identities and reproduce Table~\ref{tab:degree-two-family-counts}.
Their reports bind the input hashes and cover 421 templates, 96,036 template
states, 8,704 boundaries and 158,976 directed interfaces. The margin of two
is the finite premise for restoration; the unbounded splice and
eleven-pentagon induction remain analytic proofs.

\end{document}
